\documentclass{ltjsarticle}  % LuaLaTeX 向け
\usepackage{amsmath}
\usepackage{tikz}
\usetikzlibrary{matrix, positioning}
\usepackage{amssymb}
\usepackage{graphicx}
\usepackage{luatexja}
\usepackage{enumitem}
\usepackage{luatexja-fontspec}
\usepackage{cite}
\usepackage{hyperref}
\usepackage[a4paper,top=20mm,bottom=30mm,left=25mm,right=25mm]{geometry}
\usepackage{listings}
\usepackage{xcolor}
\usepackage{tabularx}

%%% ========================================
%%%  ここを編集してください
%%% ========================================
\date{2026年07月28日}        % 日付
\title{リスクマネジメント第14回課題}  % タイトル

\newcommand{\authorname}{安全社会基盤工学 電気システム情報コース\\学籍番号:26720493\\山口紘太郎}  % 氏名
%%% ========================================

\lstset{
  language=C,
  basicstyle=\ttfamily\footnotesize,
  keywordstyle=\color{blue},
  commentstyle=\color{gray},
  stringstyle=\color{orange},
  numbers=left,
  numberstyle=\tiny,
  numbersep=5pt,
  frame=single,
  breaklines=true,
  captionpos=b,
  showstringspaces=false,
  columns=flexible
}

\setmainjfont[
  AutoFakeBold = false
]{IPAexMincho}

\makeatletter
\renewcommand{\maketitle}{%
  \begin{flushright}
    \@date
  \end{flushright}
  \vspace{1cm}
  \begin{center}
    {\Large \@title}
  \end{center}
  \begin{flushright}
    {\large \authorname}
  \end{flushright}
  \vspace{1.5cm}
}
\makeatother

\begin{document}
\maketitle

\section{60年地震発生確率}
授業資料から，2025年の累積確率は0.39であり，2085年の累積確率は0.98となっている．
これより，2025年から2085年までの60年間に地震が発生する確率は
\[
	0.98 - 0.39 = 0.59
\]
となる．一方，2025年以降にいずれ地震が発生する確率は
\[
	1 - 0.39 = 0.61
\]いました．
である．これより，今後60年以内に地震が発生する確率は，
\[
	\frac{0.98 - 0.39}{1 - 0.39} = \frac{0.59}{0.61} \approx 0.97
\]
すなわち約97\%となる．

気づいたこととしては，30年発生確率が約75\%であったのに対し，
期間を2倍の60年にしても確率は約97\%であり，単純に2倍（約150\%）にはならないという点である．
これは確率が100\%を超えられないためで，期間を延ばすほど100\%に漸近しながら頭打ちになる．
すなわち，南海トラフ地震は長期的に見ればいつ起きてもおかしくない状況にあるといえる．

\section{講義で興味を持ったこと}
今回の講義で最も関心を持ったのは，地震発生確率という指標が存在することだ．
これまで，地震がいつ起こるかは全く予測できないものだと思っていたが，
BPTモデルのように過去の発生履歴や平均発生間隔をもとに
「今後30年以内に何\%」といった形で確率を評価できることを，
今回初めて知ることができた．

一方で，別の授業で「日本人は地震を待っている．
地震が起きて初めて動き出す」という話を聞いたことがある．
実際，大きな地震が起きてからでないと本格的な対策が進まないという体制になっていることは，
否定できないように思う．
しかし，南海トラフ地震のように，発生確率から見れば今後数十年〜100年以内には
ほぼ確実に起きることがわかっている以上，
「起きてから」ではなく「起きる前から」対策を進めることが重要だと感じた．

福島の原子力事故では，想定以上の津波など想定外の事象が発生したことによって
重大な事故となったため，南海トラフでは想定外が起きないような対策を行ってほしいと思った．

\end{document}
